\lab{13}{Ранжирование TF-IDF}
\task
Реализовать TF-IDF, разрешить частичное совпадение слов обычного
запроса, сохранить строгость явных булевых операций. Привести удачные
и неудачные примеры и сравнить поведение со строгим поиском.

\method
Вектор документа строится по нормализованным термам. Для ненулевой
частоты \(f_{t,d}\) используются веса:
\[
tf_{t,d}=1+\ln f_{t,d},\qquad
idf_t=1+\ln\frac{N+1}{df_t+1},\qquad
w_{t,d}=tf_{t,d}\,idf_t.
\]
Вес уникального терма запроса равен \(idf_t\). Повтор одного слова
в запросе не увеличивает его вес. Итоговая оценка - косинус:
\[
score(d,q)=\frac{\sum_{t\in q\cap d}(1+\ln f_{t,d})idf_t^2}
{\sqrt{\sum_{t\in d}((1+\ln f_{t,d})idf_t)^2}
\sqrt{\sum_{t\in q}idf_t^2}}.
\]
Норма документа включает весь его словарь и вычисляется заранее
при построении индекса. Если произведение норм равно нулю,
оценка принимается равной нулю. Термы отрицательных условий
используются как фильтр, а не как положительный сигнал сходства.

Кандидаты обычного запроса соединяются по ИЛИ, поэтому достаточно
частичного совпадения. При наличии явных операторов, скобок или
цитат запрос выполняется строго, а ранжирование меняет порядок
допустимых кандидатов. Нечёткий режим здесь означает ослабление
булевого условия, а не исправление опечаток. Неизвестные термы
не создают постингов и не включаются в используемый кодом набор
найденных термов запроса.

\begin{lstlisting}[language=C++]
double idf = log((double(index.header.documents) + 1)
                 / (term.documents + 1)) + 1;
query_norm += idf * idf;
// Contribution for a matched posting:
hits[j].score += (1 + log(double(cursor.positions.size)))
                 * idf * idf;
\end{lstlisting}
После нормировки результаты сортируются по убыванию score,
при равенстве - по возрастанию \(docID\). Текущая реализация
собирает всех кандидатов и полностью сортирует их, даже если
выводится только пять результатов.

\journal
Координатный индекс дополнен длинами документов и нормами.
Затем добавлен сбор уникальных положительных термов запроса,
накопление скалярного произведения, нормировка и сортировка.
Строгий и ранжированный движки используют один индекс;
в конфигурации переключается \code{search_engine=search|rank}.
Для честного сравнения сохранены обе выдачи одних 30 запросов.
Измеритель открывает основной готовый индекс и проверяет соответствие
кодека, нормализации и шага конфигурации. Он не строит дополнительный
индекс, не удаляет каталог результатов и работает без исходного корпуса.

\testing
На малом корпусе норма рассчитывается вручную и сравнивается
с допуском \(10^{-12}\). Проверяются нулевая норма пустого
документа, одинаковые веса, детерминированное разрешение равенства,
строгое И и частичное совпадение. Режим rank не должен превращать
явный \code{&&} в ИЛИ. Четыре модульных теста пройдены.
CLI-проверка дополнительно выполняет измерение дважды при недоступном
корпусе, сверяет неизменность индекса и проверяет ошибку при его отсутствии.

\results
\begin{table}[H]\centering
\begin{tabular}{lrr}\toprule
Запрос & Строгий поиск & rank (ИЛИ)\\\midrule
\code{memory allocation} & 36\,917 & 282\,185\\
\code{mutex locking} & 3\,230 & 30\,501\\
\code{waiting children} & 7\,543 & 216\,266\\
\code{file descriptors} & 16\,623 & 753\,972\\
\code{memory copies} & 85\,345 & 545\,106\\
\code{epolll} & 0 & 0\\
\bottomrule\end{tabular}
\caption{Число кандидатов в выдачах строгого и ранжированного поиска.}
\end{table}

Для запроса \code{memory allocation} первым идёт \emph{Page Allocation}
с оценкой 0,749056, затем \emph{Virtually Contiguous Memory Allocation}
с 0,539702 и \emph{C++ keyword: new} с 0,455739.
Однако первые две страницы в корпусе содержат только краткие
заголовки, а третья описывает ключевое слово C++. Поэтому все три
получили оценку релевантности 1 для запроса об инструкциях выделения
памяти в Linux и libc. Тематически подходящий заголовок ещё не означает
наличия ответа в документе. Строгий поиск начинает с
\emph{AddressSanitizer}, который также не отвечает непосредственно на запрос.

Удачный пример даёт запрос \code{malloc}: первыми выводятся
\emph{Unconstrained Allocation} и \emph{Basic Allocation} из
документации GNU libc. Для этого запроса NDCG@5 равна 1,0000.

Есть и менее удачные примеры. Для \code{memory copies} и
\code{memory overlapping} первым становится \emph{Out Of Memory
Handling}: частичное совпадение по memory не отражает намерение
найти копирование перекрывающихся областей. Для \code{printf}
первой идёт задача \emph{Number triangle with while loop in c},
а не справочная страница функции. На \code{epolll} результатов нет.
В этих случаях нужны соответственно учёт связи слов, предпочтение
API-страниц и отдельная коррекция опечаток.

Бенчмарк 30 запросов после одного прогрева и с семью повторами:
среднее 806,199 мс, медиана 264,808 мс, суммарный счётчик прочитанных
постингов 79\,385\,544. Это существенно дороже строгого
поиска, поскольку обычные запросы создают широкое объединение
кандидатов и требуют вычисления весов. Замеры разных лабораторных
не дают чистой оценки цены одной операции: у них отличаются
наборы параметров и путь исполнения.

По полностью размеченному пулу из лабораторной №6 получены
P@5 = 0,4933, DCG@5 = 7,8437, NDCG@5 = 0,4219,
ERR@5 = 0,4742. Для строгого поиска значения равны соответственно
0,1933, 2,9820, 0,1744 и 0,2283.
Сравнение оценивает весь переход от строгого поиска с порядком
по идентификатору к объединению кандидатов и TF-IDF, а не только
влияние одной формулы весов. Рост числа кандидатов сам по себе
не является улучшением качества.

\conclusions
TF-IDF выводит сначала документы с более сильным лексическим
совпадением, сохраняя строгие фильтры. Основные недостатки -
независимость термов, отсутствие смыслового понимания и высокая
цена полного ранжирования широкого множества. Улучшения:
top-k-отбор вместо полной сортировки, BM25, отдельный вес заголовка,
бонус близости слов и защита технических имён. Каждое изменение
следует проверять на одном размеченном пуле.
